\documentclass[12pt]{article}

% Packages for formatting
\usepackage[margin=1in]{geometry} % Standard 1-inch margins
\usepackage{times}                % Times New Roman font (Standard for academic SOPs)
\usepackage[utf8]{inputenc}
\usepackage{setspace}             % For line spacing
\usepackage{parskip}              % Adds space between paragraphs, removes indentation

% Set line spacing to 1.15 for readability without wasting space
\setstretch{1.0}

\begin{document}

% Header Section
\begin{center}
    {\Large \textbf{Personal Statement}} \\
    \vspace{1em}
    \textbf{Applicant:} Yang Sen (Liam) Lin \hfill \textbf{Program:} Robotics M.S.
\end{center}
\vspace{1em}

The defining cadence of my undergraduate years was set in a pool. As a varsity swimmer and City-Level Gold Medalist, I learned that physical exhaustion is temporary, but the grit required to push through it is permanent. I applied this same resilience to a significant academic barrier: pivoting from Mechanical Engineering (ME) to Computer Science (CS).

Navigating this transition was isolating; entering the world of algorithms from a background of thermodynamics initially made me feel like an outsider. However, this struggle gave me a unique vantage point. I realized my value lay not just in mastering code, but in bridging the physical constraints of hardware with the abstract logic of software. This experience clarified my purpose: to build truly robust embodied systems, I needed to unify the physical intuition of mechanics with the algorithmic reasoning of AI.

This shaped my core identity as a ``Technical Bridge.'' As Chapter Lead of the Google Developer Student Club (GDSC), I noticed non-CS majors often felt intimidated by technical jargon. Acting as a ``Translator,'' I mentored students by framing coding concepts through diverse metaphors---such as explaining API logic to a Literature major as narrative grammar. Her breakthrough reinforced my belief that cross-disciplinary friction is fuel for innovation. By orchestrating workshops for over 40 diverse students, I fostered an ecosystem where electrical engineers could fearlessly debate architecture with software theorists.

My commitment to engineering extends to ``human-centered logistics.'' Initiating a meal distribution program for the homeless in Taipei taught me that optimization metrics fail when they ignore human dignity. Managing these logistics aligned me deeply with the ethos of ``Robotics for the People,'' reminding me that a robotic system is only as robust as its ability to serve the most vulnerable.

I seek to join the University of Michigan to contribute to a community that values this breadth of perspective. The Robotics Institute is the ideal environment to refine my role as a technical bridge-builder. I plan to actively participate in the culture of the Ford Robotics Building, organizing roundtables that connect students focused on ``Reasoning'' with those focused on ``Acting.'' By sharing my experience navigating the friction between ME and CS, I hope to ensure the Michigan Robotics community remains not just intellectually elite, but profoundly inclusive.

\end{document}
